\chapter{Logic: the grammar of proof}\label{ch:logic}
\begin{readergoal}Learn the connectives, truth conditions, quantifiers, and the difference between logical implication and scientific causation.\end{readergoal}
\section{And, or, and not}
Let $P$ mean ``the animal is running'' and $Q$ mean ``theta power exceeds a specified threshold.'' The conjunction $P\land Q$ is true precisely when both claims are true. The disjunction $P\lor Q$ is false only when both are false; in mathematical usage it is \emph{inclusive}. Negation $\neg P$ reverses the truth value of $P$. The entire truth table is small enough to check directly:
\[\begin{array}{cc|ccc}P&Q&P\land Q&P\lor Q&P\Rightarrow Q\\\hline
T&T&T&T&T\\T&F&F&T&F\\F&T&F&T&T\\F&F&F&F&T\end{array}\]
The conditional $P\Rightarrow Q$ is false only when $P$ holds but $Q$ does not. This is the \emph{material implication}; it is not an assertion of a mechanism or a causal influence. Indeed, a false antecedent makes the material implication true without producing any biological event.

\section{Converse, inverse, and contrapositive}
From $P\Rightarrow Q$ one may infer the logically equivalent contrapositive $\neg Q\Rightarrow\neg P$. The converse $Q\Rightarrow P$ and inverse $\neg P\Rightarrow\neg Q$ are different claims. A low firing rate might be observed during inactivity, but inactivity need not be the only circumstance producing that rate. We may write $P\Leftrightarrow Q$ for a \emph{biconditional}, meaning that both implications hold.

De Morgan's rules supply a useful bridge to data selection: $\neg(P\land Q)\Leftrightarrow(\neg P\lor\neg Q)$ and $\neg(P\lor Q)\Leftrightarrow(\neg P\land\neg Q)$. Thus the complement of trials satisfying both inclusion criteria is the collection failing at least one criterion.

\section{Quantifiers and their order}
The phrase ``for every'' is denoted $\forall$, and ``there exists'' by $\exists$. These operations are not interchangeable: $\forall i\,\exists t\,F(i,t)$ says each neuron fires at some time, possibly a different time for every neuron. The stronger statement $\exists t\,\forall i\,F(i,t)$ claims there is one time at which every neuron fires. Neither form alone specifies simultaneity to a millisecond: that depends on how $t$ and $F$ are defined.

Negation propagates through quantifiers by the rules
\[\neg\forall x\,P(x)\ \Leftrightarrow\ \exists x\,\neg P(x),\qquad
\neg\exists x\,P(x)\ \Leftrightarrow\ \forall x\,\neg P(x).\]
For a finite dataset, verifying a universal claim requires checking all included cases; refuting it requires only one counterexample.

\section{Necessary and sufficient conditions}
If $P\Rightarrow Q$, then $P$ is \emph{sufficient} for $Q$, while $Q$ is \emph{necessary} for $P$. An example: if a matrix $A$ is invertible, its columns are linearly independent. In the square finite-dimensional case, the converse also holds; the statements are equivalent. Contrast this with ``if the signal is a sinusoid, it has a spectral peak,'' whose converse may fail for many reasons. Whenever an algorithm claims to ``detect oscillations,'' ask whether its acceptance rule is necessary, sufficient, both, or neither for the physiological interpretation.

\section{Formal proof laboratory: propositional and quantified logic}
\label{sec:lean-logic-pilot}
The five mathematical statements below are shown alongside their
Lean 4 implementations. A successful Lean build means the kernel
accepts the propositions \emph{as typed}; it does not establish
the scientific interpretation of a null hypothesis.

\begingroup
\setlength{\columnsep}{0.05\textwidth}
\noindent\begin{minipage}[t]{0.45\textwidth}
\textbf{English mathematics}

\textbf{LOGIC-001.} De Morgan's conjunction law:
for arbitrary propositions $P,Q$,
$\neg(P\land Q)\iff\neg P\lor\neg Q$.
If not both criteria hold, at least one fails. \emph{Proof:} by excluded middle, either $P$ is false, or $P$ holds, in which case $Q$ must be false. Conversely, either negated conjunct contradicts $P\land Q$.

\medskip\textbf{LOGIC-002.} De Morgan's disjunction law:
$\neg(P\lor Q)\iff\neg P\land\neg Q$.
To fail either-of-two inclusion rules is to fail both. \emph{Proof:} $\neg(P\lor Q)$ separately contradicts each disjunct. Conversely, $\neg P\land\neg Q$ refutes every proposed witness for $P\lor Q$.

\medskip\textbf{LOGIC-003.} Contraposition:
$(P\Rightarrow Q)\iff(\neg Q\Rightarrow\neg P)$.
The backward direction is classical: assuming
$\neg Q\Rightarrow\neg P$, if $P$ held but $Q$ did not,
one would obtain $\neg P$, a contradiction; by excluded middle, $Q$ must hold. The forward implication follows constructively: from $P\Rightarrow Q$ and $\neg Q$, assuming $P$ would contradict $\neg Q$.

\medskip\textbf{LOGIC-004.} Quantifier negation:
$\neg(\forall x\,R(x))\iff\exists x\,\neg R(x)$.
A counterexample disproves a universal claim. \emph{Proof:} a failing witness contradicts the universal assertion. Conversely, if no witness failed, classical double-negation elimination would give $R(x)$ for every $x$, contradicting $\neg\forall x\,R(x)$.

\medskip\textbf{LOGIC-005.} Existential negation:
$\neg(\exists x\,R(x))\iff\forall x\,\neg R(x)$. \emph{Proof:} if a witness existed, it would contradict the assumption that none exists. Conversely, if every candidate fails, an assumed witness immediately yields a contradiction.

\medskip\textbf{Proof strategy.} The statements about negating
conjunction, contraposition and negating a universal
quantifier use classical reasoning. The Lean proof
invokes Mathlib tactics that elaborate these elementary
steps; the file prints their axiom dependencies.
\end{minipage}\hfill
\begin{minipage}[t]{0.50\textwidth}
\textbf{Lean 4 / Mathlib}

\lstinputlisting[language=LeanFour,basicstyle=\ttfamily\tiny,
breaklines=true,columns=fullflexible,keepspaces=true,
frame=single]{../../formal/lean/Con2physFormal/Logic.lean}

\textit{Correspondence and provenance:}
each Lean declaration has the same quantification
over propositions or predicates as the adjacent English
statement. Mathlib provides the underlying logic and tactics.
See the pinned dependency in
\texttt{\detokenize{formal/lean/lakefile.lean}}.
\end{minipage}
\endgroup


\begin{workedexample}[A null hypothesis in plain language]
Suppose $H_0$ states that the temporal positions of spikes inside specified jitter windows carry no additional information beyond their counts. Rejecting $H_0$ is evidence against this particular conditional model. It does \emph{not} establish that a synapse connects the neurons, or that all alternative null models would also be rejected. Compare \cite{Amarasingham2012}.
\end{workedexample}
% --- expanded edition (October 2026) ---
\section{Tautologies, contradictions, and valid arguments}
A compound formula is a \emph{tautology} if it is true under every assignment of truth values to its atomic letters, and a \emph{contradiction} if it is false under every assignment. The law of excluded middle $P\lor\neg P$ is a tautology; $P\land\neg P$ is a contradiction. Two formulas $\varphi$ and $\psi$ are \emph{logically equivalent}, written $\varphi\equiv\psi$, when $\varphi\Leftrightarrow\psi$ is a tautology.

An \emph{argument} consists of premises $\varphi_1,\ldots,\varphi_k$ and a conclusion $\psi$. It is \emph{valid} when every assignment making all premises true also makes $\psi$ true; equivalently, when $(\varphi_1\land\cdots\land\varphi_k)\Rightarrow\psi$ is a tautology. Validity concerns form, not content: a valid argument with a false premise can have a false conclusion, and an invalid argument can happen to have a true conclusion.

\begin{proposition}[Four standard inference rules]
The following argument forms are valid:
\begin{enumerate}[leftmargin=2em]
\item \emph{modus ponens}: from $P$ and $P\Rightarrow Q$, infer $Q$;
\item \emph{modus tollens}: from $\neg Q$ and $P\Rightarrow Q$, infer $\neg P$;
\item \emph{hypothetical syllogism}: from $P\Rightarrow Q$ and $Q\Rightarrow R$, infer $P\Rightarrow R$;
\item \emph{disjunctive syllogism}: from $P\lor Q$ and $\neg P$, infer $Q$.
\end{enumerate}
\end{proposition}
\begin{proof}
We check (2); the others are similar. Suppose $\neg Q$ and $P\Rightarrow Q$ are both true. If $P$ were true, then by the truth table of the conditional $Q$ would be true, contradicting $\neg Q$. Hence $P$ is false, i.e.\ $\neg P$ is true. For (3), suppose both premises are true and $P$ is true; then $Q$ is true by the first premise and $R$ by the second, so $P\Rightarrow R$ cannot be falsified. For (4), if $P\lor Q$ is true and $P$ is false, the only remaining way for the disjunction to hold is that $Q$ is true.
\end{proof}

Two \emph{invalid} forms are so common that they have names. \emph{Affirming the consequent} infers $P$ from $Q$ and $P\Rightarrow Q$; \emph{denying the antecedent} infers $\neg Q$ from $\neg P$ and $P\Rightarrow Q$. In both cases the assignment $P=F$, $Q=T$ makes the premises true and the conclusion false. A frequent scientific instance: ``if the region generates gamma oscillations, the LFP shows a gamma peak; the LFP shows a gamma peak; therefore the region generates gamma.'' Volume conduction from another region, or a harmonic of a lower rhythm, also produces the peak.

\section{Useful equivalences}
The following equivalences can each be verified with a four- or eight-row truth table and will be used without comment later.
\[
\begin{array}{rcl@{\qquad}l}
P\Rightarrow Q&\equiv&\neg P\lor Q & \text{(material implication)}\\
P\Rightarrow Q&\equiv&\neg Q\Rightarrow\neg P & \text{(contraposition)}\\
\neg(P\Rightarrow Q)&\equiv&P\land\neg Q & \text{(negated conditional)}\\
P\land(Q\lor R)&\equiv&(P\land Q)\lor(P\land R) & \text{(distributivity)}\\
(P\land Q)\Rightarrow R&\equiv&P\Rightarrow(Q\Rightarrow R) & \text{(exportation)}\\
P\Leftrightarrow Q&\equiv&(P\Rightarrow Q)\land(Q\Rightarrow P) & \text{(biconditional)}
\end{array}
\]
The negated conditional is the most practically useful: to show that ``$P$ implies $Q$'' fails, one must exhibit a case where $P$ holds and $Q$ does not. A case where $P$ fails is irrelevant.

\begin{proposition}
For any propositions, $\neg(P\Rightarrow Q)\equiv P\land\neg Q$.
\end{proposition}
\begin{proof}
By material implication, $\neg(P\Rightarrow Q)\equiv\neg(\neg P\lor Q)$. By De Morgan's rule this is $\neg\neg P\land\neg Q$, and double negation gives $P\land\neg Q$.
\end{proof}

\section{Nested quantifiers and their negations}
Many definitions in analysis have three or more alternating quantifiers. The definition of a limit, met in Chapter~\ref{ch:limits}, reads
\[
\forall\varepsilon>0\ \exists N\ \forall n\ge N:\ |x_n-L|<\varepsilon.
\]
Its negation is obtained mechanically by moving $\neg$ inward one quantifier at a time, swapping $\forall$ and $\exists$:
\[
\exists\varepsilon>0\ \forall N\ \exists n\ge N:\ |x_n-L|\ge\varepsilon.
\]
In words: there is a fixed tolerance that infinitely late terms keep violating. Note that the bounded quantifiers $\forall\varepsilon>0$ and $\exists n\ge N$ keep their bounds under negation; only the quantifier type and the inner predicate change.

The \emph{uniqueness quantifier} $\exists!x\,P(x)$ (``there is exactly one $x$ with $P$'') abbreviates
\[
\exists x\,\bigl(P(x)\land\forall y\,(P(y)\Rightarrow y=x)\bigr).
\]
Existence and uniqueness are proved separately: one exhibits an $x$, then shows that any two such objects coincide.

\begin{workedexample}[Negating a selectivity criterion]
Suppose a unit is called ``direction selective'' if there is a direction $d^\ast$ such that for every other tested direction $d$, the mean rate at $d^\ast$ exceeds the rate at $d$ by at least 5 spikes/s. Symbolically, $\exists d^\ast\,\forall d\ne d^\ast:\ r(d^\ast)-r(d)\ge5$. Its negation, ``not direction selective'', is $\forall d^\ast\,\exists d\ne d^\ast:\ r(d^\ast)-r(d)<5$: every candidate preferred direction has at least one competitor within 5 spikes/s. This is a much weaker statement than ``the unit responds equally to all directions,'' and the two must not be conflated in a results section.
\end{workedexample}

\section{Logic and uncertainty}
Classical logic assigns each proposition exactly one of two truth values. Empirical claims usually come with uncertainty, and it is tempting to treat a $p$-value or posterior probability as a ``degree of truth.'' This is a category mistake. A probability is assigned by a model to an event; the proposition ``this event occurred'' remains true or false. Probability theory (Part~V) extends logic by attaching numbers to sets of outcomes, but every probabilistic statement is itself a classical proposition about those numbers, and proofs about probabilities use exactly the rules of this chapter.
\begin{exercises}
\item Construct a complete truth table for $(P\Rightarrow Q)\Leftrightarrow(\neg Q\Rightarrow\neg P)$.
\item Negate ``for every frequency there exists a channel with nonzero power.''
\item Give a counterexample to the inference from $P\Rightarrow Q$ and $Q$ to $P$.
\item Explain why statistical dependence is not equivalent to causal influence.
\item Use a truth table to show that $(P\Rightarrow Q)\land(Q\Rightarrow R)\Rightarrow(P\Rightarrow R)$ is a tautology.
\item Show that affirming the consequent is invalid by exhibiting a falsifying truth assignment.
\item Negate $\forall\varepsilon>0\ \exists\delta>0\ \forall x\ (|x-a|<\delta\Rightarrow|f(x)-f(a)|<\varepsilon)$ and state the result in words.
\item Write ``exactly one channel exceeds the artifact threshold'' using $\exists!$ and then without it.
\item Prove the exportation law $(P\land Q)\Rightarrow R\equiv P\Rightarrow(Q\Rightarrow R)$.
\end{exercises}
